\documentclass[11pt]{article}
\usepackage[margin=1in]{geometry}
\usepackage{amsmath,amssymb,amsthm}
\usepackage{hyperref}
\usepackage{xurl}

\newtheorem{theorem}{Theorem}
\newtheorem{lemma}[theorem]{Lemma}
\newtheorem{corollary}[theorem]{Corollary}
\theoremstyle{definition}
\newtheorem{definition}[theorem]{Definition}
\theoremstyle{remark}
\newtheorem{remark}[theorem]{Remark}

\let\oldus\_
\renewcommand{\_}{\oldus\allowbreak}
\newcommand{\Cl}{\operatorname{Cl}}
\newcommand{\Idem}{\operatorname{Idem}}
\newcommand{\ZD}{\operatorname{ZD}}

\title{A proof of Conjecture 00000007630\\[2pt]
\large Clifford algebras of signatures $(0,1,0)$ and $(0,0,1)$: same idempotents, different zero divisors}
\author{Jackmeson1}
\date{}

\begin{document}
\sloppy
\maketitle

\section{The conjecture}

Conjecture 00000007630 reads (English version; the Chinese version agrees, and it says
``explicit signature pair'' for the last clause):

\begin{quote}
\emph{Definition:} Idempotent spectra and zero-divisor spectra are two layers.
\emph{Conjecture:} There exist two Clifford algebras with the same idempotent spectrum but different
zero-divisor spectra, and the separation is realized by an explicit signature pair with the same
idempotents but different zero divisors.
\end{quote}

This is an existence claim. We exhibit an explicit pair of real quadratic forms with signatures
$(0,1,0)$ and $(0,0,1)$. We prove that their Clifford algebras have the same idempotents and
different zero divisors.

\section{Reading and definitions}

\begin{definition}
For a real quadratic form $Q$ on $V$, the Clifford algebra $\Cl(Q)$ is the quotient of the tensor
algebra $T(V)$ by the relations $v\otimes v=Q(v)\cdot 1$. In Lean this is Mathlib's
\texttt{CliffordAlgebra Q}, with generator map $\iota\colon V\to\Cl(Q)$. Degenerate forms are
allowed, as in Mathlib; the Clifford algebra of the zero form is the exterior algebra. The
\emph{signature} of $Q$ is its Sylvester signature $(p,q,r)$: the numbers of positive, negative
and null directions.
\end{definition}

\begin{definition}
For an algebra $A$, the idempotent spectrum is $\Idem(A)=\{x: x^2=x\}$. The zero-divisor spectrum
is $\ZD(A)=\{x:\exists y\neq0,\ xy=0 \text{ or } yx=0\}$. With this convention $0\in\ZD(A)$; nonzero
zero divisors are reported separately.
\end{definition}

\paragraph{The explicit signature pair.} On $V=\mathbb{R}$:
\begin{itemize}
\item $Q_1(t)=-t^2$ (Mathlib's \texttt{CliffordAlgebraComplex.Q}). It is negative definite, so its
signature is $(0,1,0)$. Mathlib's \texttt{CliffordAlgebraComplex.equiv} gives
$\Cl(Q_1)\cong\mathbb{C}$.
\item $Q_0(t)=0$. It is totally null, so its signature is $(0,0,1)$. Mathlib's
\texttt{CliffordAlgebraDualNumber.equiv} gives $\Cl(Q_0)\cong\mathbb{R}[\varepsilon]$ (the dual
numbers, with $\varepsilon^2=0$).
\end{itemize}

\section{The proof}

\begin{lemma}\label{lem:dual}
In the dual numbers $\mathbb{R}[\varepsilon]$, the only idempotents are $0$ and $1$.
\end{lemma}
\begin{proof}
Write $z=a+b\varepsilon$. Then $z^2=a^2+2ab\,\varepsilon$, so $z^2=z$ gives $a^2=a$ and $2ab=b$.
If $a=0$ then $b=0$; if $a=1$ then $2b=b$, so $b=0$.
\end{proof}

\begin{theorem}[same idempotents]\label{thm:idem}
$\Idem(\Cl(Q_1))=\{0,1\}$ and $\Idem(\Cl(Q_0))=\{0,1\}$. In each algebra $0\neq1$, so each has
exactly two idempotents.
\end{theorem}
\begin{proof}
$\Cl(Q_1)\cong\mathbb{C}$ has no zero divisors, and in such a ring $x^2=x$ forces $x\in\{0,1\}$.
For $\Cl(Q_0)$, transport Lemma~\ref{lem:dual} along the algebra isomorphism.
\end{proof}

\begin{theorem}[different zero divisors]\label{thm:zd}
$\ZD(\Cl(Q_1))=\{0\}$. In contrast, $\ZD(\Cl(Q_0))=\iota(\mathbb{R})=\{\iota(r):r\in\mathbb{R}\}$,
and $e=\iota(1)$ satisfies $e\neq0$ and $e^2=0$. Consequently $\Cl(Q_1)\not\cong\Cl(Q_0)$ as
$\mathbb{R}$-algebras.
\end{theorem}
\begin{proof}
The first claim holds because $\mathbb{C}$ has no zero divisors. In $\Cl(Q_0)$ we have
$\iota(r)\iota(s)=rs\,\iota(1)^2=rs\,Q_0(1)=0$, and $\iota(1)\mapsto\varepsilon\neq0$. Conversely, if
the image $a+b\varepsilon$ of $x$ has $a\neq0$, then it is a unit, so $x$ is not a zero divisor.
If $a=0$, then $x=\iota(b)$. An algebra isomorphism would carry the nonzero zero divisor $e$ to a
nonzero zero divisor of $\Cl(Q_1)$, which does not exist.
\end{proof}

Theorems~\ref{thm:idem} and~\ref{thm:zd} give the conjecture. \qed

\begin{remark}[scope and conventions]
(1) ``The same idempotent spectrum'' is proved in three senses: equal idempotent sets $\{0,1\}$,
equal cardinality ($2$), and, implicitly, the obvious bijection $0\mapsto0$, $1\mapsto1$. (2)
``Different zero-divisor spectra'' holds whether or not $0$ counts as a zero divisor: $\Cl(Q_1)$
has no nonzero zero divisor, while $\Cl(Q_0)$ has a whole line of them. (3) The second form is
degenerate. Among \emph{nondegenerate} real Clifford algebras, ``only the trivial idempotents''
already forces a division algebra ($\mathbb{R},\mathbb{C},\mathbb{H}$), since such algebras are
semisimple (Wedderburn theory). This standard fact is context only and is not used in the proof. So
under the reading ``idempotent set $=\{0,1\}$'', a separating pair must involve a degenerate
signature. This is why the null direction $r=1$ appears.
\end{remark}

\section{Lean formalization}

File \texttt{lean/Conjecture7630/Basic.lean} (Lean 4.33.1, Mathlib \texttt{v4.33.1}). It defines
\texttt{idempotents}, \texttt{zeroDivisors}, \texttt{Q1} $:=$ \texttt{CliffordAlgebraComplex.Q},
\texttt{Q0} $:= 0$, \texttt{Cl1}, \texttt{Cl0}, and \texttt{gen} $:=\iota(1)$. The theorems are:
\begin{itemize}
\item \texttt{Q1\_negDef}, \texttt{Q0\_null}: the signatures.
\item \texttt{dual\_idem}, \texttt{idempotents\_Cl1}, \texttt{idempotents\_Cl0}: Lemma~\ref{lem:dual}
and Theorem~\ref{thm:idem}.
\item \texttt{zeroDivisors\_Cl1}, \texttt{zeroDivisors\_Cl0}, \texttt{gen\_ne\_zero},
\texttt{gen\_mul\_gen}, \texttt{not\_isomorphic}: Theorem~\ref{thm:zd}.
\item \texttt{separation}: the main theorem. It bundles both signatures, both idempotent sets
$=\{0,1\}$ with \texttt{ncard} $2$, both zero-divisor sets, the nonzero square-zero element, and
non-isomorphism.
\end{itemize}
\texttt{lake build} succeeds. \texttt{\#print axioms} for \texttt{separation},
\texttt{not\_isomorphic} and \texttt{zeroDivisors\_Cl0} reports only
\texttt{[propext, Classical.choice, Quot.sound]}. The code contains no \texttt{sorry} and no
\texttt{native\_decide}.

\end{document}
